\documentclass[11pt]{article}
\usepackage[margin=1in]{geometry}
\usepackage[T1]{fontenc}
\usepackage{amsmath,amssymb,graphicx,booktabs,xcolor}
\usepackage{placeins}
\usepackage[hidelinks]{hyperref}
\setlength{\emergencystretch}{2em}
\newcommand{\todo}[1]{\par\noindent\textcolor{red}{\textbf{TODO:} #1}\par}
\title{Noise-aware optical architecture search for constrained relative-phase imaging}
\author{Author list and affiliations to be agreed}
\date{Working manuscript --- results pending}
\begin{document}
\maketitle
\begin{abstract}
Optical layouts that produce high normalized phase contrast need not preserve
enough photons for accurate measurements. We investigate joint discrete and
continuous design of coherent imaging systems under a fixed incident illumination
budget and bounded hardware geometry. The framework selects optical component
families at sequential stations and optimizes their continuous parameters using
differentiable scalar wave propagation. A shared affine phase readout permits
analytic evaluation of expected phase error under Poisson shot noise and Gaussian
read noise. Separate training, validation and held-out datasets support comparisons
with random architecture proposals, restricted libraries, fixed layouts and
freeform phase masks. \textbf{This draft specifies the implemented simulation
method; comparative results and physical validation remain to be obtained.}
\end{abstract}

\section{Introduction}
Optical design increasingly combines physical simulation with task-specific
optimization. Existing work includes evolutionary topology and continuous
parameter optimization in photonic systems~\cite{aesop}, automated discovery of
microscopy arrangements~\cite{xlumina}, discrete--continuous compound lens
design~\cite{teh}, reconfigurable diffractive systems and lens-component
deletion~\cite{chord,zhang}, diffractive quantitative phase imaging~\cite{mengu}, and
experimentally demonstrated phase microscopy with learned metasurface
optics~\cite{neuralphase}. These establish substantial prior art for the broad
ingredients of this study.

Our question is whether categorical selection among bounded conventional
free-space components improves relative-phase estimation under explicit photon
and geometric constraints. We use Chromatix~\cite{chromatix} for differentiable
wave propagation and compare alternative design strategies under the same
measurement model. The intended contribution is an evaluated constrained imaging
study and reproducible search framework, conditional on the evidence below.
\todo{Update the full novelty matrix near submission. Do not claim the first
optical architecture search, differentiable optical design, or optical phase
recovery. State the demonstrated contribution after experiments.}

\section{Methods}
\subsection{Object, optics and sensor}
For a phase map $\phi$ defined relative to a known zero-phase background, the
illuminated object field is $u_0=A\exp(i\phi)$. The current model assumes unit
object transmission within a finite square illumination region, coherent
monochromatic paraxial scalar propagation, and no polarization or multiple
scattering. Default phantoms have peak relative phase at most $0.5$ rad.
An intensity measurement cannot determine a global phase offset.

A candidate comprises $K$ sequential stations, each containing a categorical
element or identity, and $K+1$ free-space gaps. The gaps satisfy
\begin{equation}
 z_j=z_{\min}+[L-(K+1)z_{\min}]\,
 \frac{\exp r_j}{\sum_k\exp r_k},\qquad \sum_j z_j=L.
\end{equation}
The library includes thin lenses, hard and soft pupils, axicons, gratings and a
circular phase disk. A separate baseline uses bounded freeform phase masks.
Focal lengths and widths are bounded by sigmoid transforms; phase heights use a
scaled hyperbolic tangent. Installed components have finite apertures and explicit
transmission. An upper bound limits their count. These constraints do not yet
represent a discrete vendor catalog or a complete aberration model.

The field is padded once and cropped at the detector. FFT propagation retains
periodic outer boundaries. Grid and padding convergence are assessed at fixed
physical parameters. The detector currently uses point samples rather than
finite-area pixel integration. Optical intensity $I_i$ is expressed relative to
the fixed incident intensity, never normalized by the output image or its truth.

\subsection{Noise-aware phase objective}
For incident photons per input pixel $D$, detector quantum efficiency $q$, and
read standard deviation $s$ in electrons, the modeled measurement is
\begin{equation}
 Y_i=P_i+\epsilon_i,\quad P_i\sim\operatorname{Poisson}(DqI_i),\quad
 \epsilon_i\sim\mathcal N(0,s^2).
\end{equation}
Noise samples are independent across pixels. With $J_i=Y_i/(Dq)$ and affine
readout $\widehat\phi_i=aJ_i+b$, the exact expected squared error is
\begin{equation}\label{eq:risk}
 \mathbb E[(\widehat\phi_i-\phi_i)^2]
 =(aI_i+b-\phi_i)^2+a^2\left[\frac{I_i}{Dq}+\frac{s^2}{(Dq)^2}\right].
\end{equation}
This follows from the mean and variance of the specified measurement and requires
no Gaussian approximation to the Poisson distribution. Saturation, dark current,
ADC quantization and detector correlations are not currently included.

For each candidate, readout coefficients are fitted only on the training ROI:
\begin{equation}
 a=\frac{\operatorname{Cov}(I,\phi)}{\operatorname{Var}(I)+
       \mathbb E[v(I)]+\varepsilon},\qquad
 b=\mathbb E[\phi]-a\mathbb E[I],
\end{equation}
where $v(I)=I/(Dq)+s^2/(Dq)^2$ and $\varepsilon$ is a small conditioning ridge.
An upright or inverted sensor-coordinate convention is selected by training
risk with the same options for all methods. Calibration is fixed during
validation, test, and dose/defocus sweeps. The differentiable objective profiles
through this training-only fit; orientation is a piecewise discrete choice.
The final objective averages Eq.~\ref{eq:risk} over the ROI and optional defocus
ensemble, with an optional penalty per installed component.

\subsection{Search and comparisons}
The memetic algorithm alternates local Adam updates with validation selection
and categorical/continuous mutation. Sampled hard-pupil and phase-disk widths
are mutated without gradient updates. Adam moments reset at each generation;
elite optical parameters are retained. A separate archive preserves the best
evaluated design before breeding.

The main controls are random architecture proposals with the same local update
budget, a restricted element library, an optimized fixed lens--phase-disk--lens
sequence, and optimized freeform masks. A noise-free objective ablation omits
noise from training, readout fitting and selection, but uses the same noisy test.
Defocus-only optics and a constant training-mean predictor test whether useful
phase information is extracted. The six optimized methods match population,
generations and gradient steps; parameter counts and actual wall times differ.
The fixed/freeform controls are generic implementations, not reproductions of
published networks or commercial microscopes.

\subsection{Data and evaluation protocol}
Default data are independently seeded phase phantoms; the shape-shift set uses
sharp ellipses. Custom calibrated phase maps can be imported with group-disjoint
acquisitions. Training fits optical parameters and readout, validation selects
the design, and the saved selection is checksummed before the separate test
command. The primary endpoint is held-out phase RMSE at training dose and nominal
sensor position. Secondary endpoints include dose response, defocus, shape shift,
throughput and exact sampled-camera error. The readout is not refitted for these
conditions. Each evaluation uses one modeled exposure per specimen.

RMSE is calculated per independent search seed before averaging. The generated
descriptive bootstrap intervals resample seeds conditional on the fixed test
specimens; pixels and repeated noise draws are not independent biological samples.
\todo{Specify the final datasets, physical sampling, hyperparameter budget,
effect-size target, seed count, convergence tolerances, hardware calibration,
and specimen-level uncertainty analysis before unblinding the main study.}

\section{Results}
\todo{Report all prespecified methods, failed runs and budgets. Include null or
negative findings. Smoke tests establish execution only.}
\IfFileExists{generated/results_table.tex}{%
\begin{table}[ht]\centering\small
\input{generated/results_table.tex}
\caption{Automatically generated nominal-condition simulation results. Confirm
study identity, final protocol and adequate replication before publication.}
\end{table}}{\todo{Insert the table generated from the frozen evaluated study.}}
\IfFileExists{generated/dose_accuracy.pdf}{%
\begin{figure}[ht]\centering
\includegraphics[width=.85\linewidth]{generated/dose_accuracy.pdf}
\caption{Simulated dose response of fixed trained designs and readouts. Bands,
when present, describe variation across search seeds on a fixed test set.}
\end{figure}}{\todo{Insert the generated dose--accuracy figure after validation.}}
\todo{Add fixed-design convergence, architecture/complexity analysis, strong
published-method controls and measured optical validation.}

\FloatBarrier
\section{Discussion}
Any observed gain must be separated from numerical artifacts, unequal tuning,
readout capacity, and assumptions about photon or component efficiency. The
current model is restricted to weak relative phase, simple calibration, a bounded
sequential path and ideal scalar components. Synthetic shape generalization does
not establish biological generalization. A convincing experimental claim needs
measured detector/optics calibration, phase standards, independent specimens and
appropriate reference-exposure accounting. Global optimality and uniqueness of a
recovered optical layout are not implied by heuristic search.
\todo{Discuss demonstrated advantages and failure modes only after results exist.}

\section*{Code and data availability}
Code is available in the
\href{https://github.com/Le0nZim/optical-architecture-search}{Optical Architecture Search repository}.
Runs export configurations, physical prescriptions, calibration, code hashes,
dependencies, dataset hashes and per-specimen metrics.
\todo{Archive the final release and permitted datasets with permanent identifiers;
specify the exact release used. Confirm author contributions and funding.}
\bibliographystyle{unsrt}
\bibliography{references}
\end{document}
